% Compilar dos veces con pdflatex desde este directorio.
\documentclass[aspectratio=169,10pt]{beamer}
\usepackage[utf8]{inputenc}
\usepackage[T1]{fontenc}
\usepackage[spanish,es-nodecimaldot]{babel}
\usepackage{amsmath,amssymb,booktabs,tabularx,array,tikz}
\usetikzlibrary{arrows.meta,positioning,calc,shapes.geometric,fit,backgrounds}

\definecolor{UAPNavy}{HTML}{17324D}
\definecolor{UAPTeal}{HTML}{007F82}
\definecolor{UAPGold}{HTML}{E2A33A}
\definecolor{UAPLight}{HTML}{F3F6F8}
\definecolor{UAPGray}{HTML}{536471}
\definecolor{UAPRed}{HTML}{A94343}
\definecolor{OuterBlue}{HTML}{2878A8}
\definecolor{InnerGreen}{HTML}{358A68}

\tikzset{
  flow/.style={-{Stealth[length=2.1mm]},thick,draw=UAPNavy},
  data/.style={trapezium,trapezium left angle=70,trapezium right angle=110,draw=UAPGold,fill=UAPGold!14,rounded corners=1pt,align=center,minimum height=8mm,text=UAPNavy},
  process/.style={draw=UAPTeal,fill=UAPTeal!11,rounded corners=3pt,very thick,align=center,minimum height=9mm,text=UAPNavy},
  external/.style={draw=OuterBlue,fill=OuterBlue!8,rounded corners=3pt,very thick,align=center,minimum height=9mm,text=UAPNavy},
  internal/.style={draw=InnerGreen,fill=InnerGreen!9,rounded corners=3pt,very thick,align=center,minimum height=9mm,text=UAPNavy},
  decision/.style={diamond,aspect=1.7,draw=UAPNavy,fill=UAPNavy!7,very thick,align=center,inner sep=1.5pt,text=UAPNavy},
  optional/.style={draw=UAPGray,dashed,fill=UAPGray!7,rounded corners=3pt,align=center,minimum height=9mm,text=UAPNavy},
  outerbox/.style={draw=OuterBlue,very thick,rounded corners=4pt,inner sep=5pt},
  innerbox/.style={draw=InnerGreen,very thick,rounded corners=4pt,inner sep=4pt}
}

\usetheme{Boadilla}
\usecolortheme{default}
\setbeamercolor{normal text}{fg=UAPNavy,bg=white}
\setbeamercolor{structure}{fg=UAPTeal}
\setbeamercolor{frametitle}{fg=white,bg=UAPNavy}
\setbeamercolor{title}{fg=white}
\setbeamercolor{subtitle}{fg=UAPGold}
\setbeamercolor{block title}{fg=white,bg=UAPTeal}
\setbeamercolor{block body}{fg=UAPNavy,bg=UAPLight}
\setbeamercolor{alertblock title}{fg=white,bg=UAPRed}
\setbeamercolor{alertblock body}{fg=UAPNavy,bg=UAPRed!7}
\setbeamerfont{frametitle}{series=\bfseries}
\setbeamertemplate{navigation symbols}{}
\setbeamertemplate{itemize item}{\color{UAPGold}$\blacktriangleright$}
\setbeamertemplate{itemize subitem}{\color{UAPTeal}$\bullet$}
\setbeamertemplate{footline}{%
  \leavevmode\hbox{%
    \begin{beamercolorbox}[wd=.78\paperwidth,ht=2.8ex,dp=1.2ex,leftskip=1em]{author in head/foot}%
      \color{white}\insertshorttitle
    \end{beamercolorbox}%
    \begin{beamercolorbox}[wd=.22\paperwidth,ht=2.8ex,dp=1.2ex,rightskip=1em plus1fil]{date in head/foot}%
      \color{UAPGray}\hfill Semana 7 \enspace | \enspace \insertframenumber/\inserttotalframenumber
    \end{beamercolorbox}}}

\newcommand{\concepto}[1]{\textcolor{UAPTeal}{\textbf{#1}}}
\newcommand{\idea}[1]{\begin{block}{Idea central}#1\end{block}}
\title[Estructura base de un proyecto de ML]{Estructura base de un proyecto de Machine Learning}
\subtitle{Selección interna, evaluación externa y ajuste definitivo}
\author{Curso de Inteligencia Artificial}
\institute{Semana 7 · Presentación complementaria}
\date{}

\begin{document}

% 1
\begin{frame}[plain]
  \begin{beamercolorbox}[wd=\paperwidth,ht=\paperheight,sep=2em]{frametitle}
    \vfill
    {\Huge\bfseries\inserttitle\par}\vspace{.45em}
    {\color{UAPGold}\Large\insertsubtitle\par}\vspace{1.2em}
    \begin{tikzpicture}[node distance=5mm,scale=.75,transform shape]
      \node[data] (p) {Protocolo};
      \node[internal,right=of p] (s) {Selección interna};
      \node[external,right=of s] (e) {Evaluación externa};
      \node[process,right=of e] (f) {Ajuste final};
      \draw[flow] (p)--(s); \draw[flow] (s)--(e); \draw[flow] (e)--(f);
    \end{tikzpicture}\par
    \vfill {\color{UAPGold}\large\insertauthor\par}{\color{white}\insertinstitute\par}
  \end{beamercolorbox}
\end{frame}

% 2
\begin{frame}{El objeto evaluado es el procedimiento}
  Un proyecto de ML no evalúa solamente un algoritmo. Evalúa una composición:
  \[
    \text{datos}\rightarrow\text{transformaciones}\rightarrow\text{representación}
    \rightarrow\text{modelo}\rightarrow\text{regla de decisión}.
  \]
  \begin{itemize}
    \item Imputación, categorías, escalas y componentes pueden aprender parámetros.
    \item Selección de variables, hiperparámetros y umbral también usan datos.
    \item Toda decisión adaptativa forma parte del procedimiento que debe aislarse.
  \end{itemize}
  \idea{La unidad evaluada es el pipeline completo, no el estimador aislado.}
\end{frame}

% 3
\begin{frame}{Protocolo antes de ajustar}
  \small
  \begin{columns}[T]
    \column{.49\textwidth}
    \begin{block}{Formulación}
      \begin{itemize}
        \item problema, población y unidad;
        \item instante y horizonte;
        \item información disponible;
        \item salida y uso posterior.
      \end{itemize}
    \end{block}
    \column{.49\textwidth}
    \begin{block}{Evaluación}
      \begin{itemize}
        \item alternativas y configuraciones;
        \item métrica y regla de selección;
        \item partición por unidad, grupo o tiempo;
        \item prueba, semillas y presupuesto.
      \end{itemize}
    \end{block}
  \end{columns}
  \begin{alertblock}{Alcance}La anidación separa selección y evaluación, pero no corrige una formulación o partición inadecuada.\end{alertblock}
\end{frame}

% 4
\begin{frame}{Diagrama general de validación anidada}
  \centering
  \begin{tikzpicture}[scale=.92,transform shape,x=1cm,y=1cm,font=\footnotesize]
    % Una fila de ajuste y otra de evaluación; el retorno tiene su propio carril.
    \node[data,text width=2cm] (trainout) at (1.2,0) {Entrenamiento\\externo};
    \node[internal,text width=1.65cm] (fitin) at (4.0,0) {Ajustar\\candidatos};
    \node[internal,text width=1.65cm] (valin) at (6.6,0) {Validar y agregar\\folds internos};
    \node[internal,text width=1.65cm] (choose) at (9.2,0) {Elegir\\configuración};
    \node[external,text width=2cm] (refit) at (12.0,0) {Reajustar con\\entrenamiento externo};

    \node[external,text width=1.8cm] (predict) at (12.0,-1.8) {Predecir y guardar\\fuera de muestra};
    \node[decision,text width=1.55cm] (more) at (7.0,-1.8) {¿Más folds\\externos?};
    \node[process,text width=1.9cm] (aggregate) at (1.5,-1.8) {Agregar predicciones\\y calcular métricas};
    \node[data,text width=1.65cm] (valout) at (12.0,-3.15) {Validación\\externa};

    \draw[flow] (trainout)--(fitin);
    \draw[flow] (fitin)--(valin);
    \draw[flow] (valin)--(choose);
    \draw[flow] (choose)--(refit);
    \draw[flow] (refit)--(predict);
    \draw[flow] (valout)--(predict);
    \draw[flow] (predict)--(more);
    \draw[flow] (more)--node[above] {no: finalizar} (aggregate);
    \draw[flow] (more.south) -- (7,-3.15)
      -- node[below] {sí: siguiente partición externa} (-.25,-3.15)
      |- (trainout.west);

    \node[text=InnerGreen,font=\footnotesize\bfseries] at (6.6,.9) {Selección interna: repetir en todos los folds internos};
    \node[text=OuterBlue,font=\footnotesize\bfseries,anchor=west] at (-.55,1.5) {Evaluación externa: cada partición de los datos de desarrollo};
    \begin{scope}[on background layer]
      \node[innerbox,fit=(fitin)(valin)(choose),inner sep=8pt] {};
      \draw[outerbox] (-.65,-3.65) rectangle (13.25,1.22);
    \end{scope}
  \end{tikzpicture}

  \vspace{1mm}
  \scriptsize Los dos conjuntos externos se renuevan en cada partición.\\
  La prueba opcional se muestra en la diapositiva 9.
\end{frame}

% 5
\begin{frame}{Ciclo externo: estimar el procedimiento}
  \begin{center}
    \begin{tikzpicture}[node distance=6mm,scale=.75,transform shape]
      \node[data] (d) {Desarrollo};
      \node[external,right=of d] (s) {Partición externa};
      \node[data,right=of s,yshift=8mm] (tr) {Entrenamiento\\externo};
      \node[data,right=of s,yshift=-8mm] (va) {Validación\\externa};
      \node[internal,right=8mm of tr] (sel) {Selección\\interna};
      \node[external,right=8mm of sel] (fit) {Reajuste};
      \node[external,right=8mm of va] (pred) {Predicción\\externa};
      \draw[flow] (d)--(s); \draw[flow] (s)--(tr); \draw[flow] (s)--(va);
      \draw[flow] (tr)--(sel); \draw[flow] (sel)--(fit); \draw[flow] (fit)--(pred); \draw[flow] (va)--(pred);
    \end{tikzpicture}
  \end{center}
  \small
  \begin{enumerate}
    \item toda selección ocurre dentro del entrenamiento externo;
    \item la configuración elegida se reajusta con ese entrenamiento completo;
    \item la validación externa se predice una sola vez;
    \item se conservan predicciones fuera de muestra.
  \end{enumerate}
  \begin{block}{Terminología}“Externa” describe su posición respecto de la selección; no implica necesariamente otra población.\end{block}
\end{frame}

% 6
\begin{frame}{Ciclo interno: elegir la configuración}
  \begin{center}
    \begin{tikzpicture}[node distance=5mm,scale=.72,transform shape]
      \node[data] (tr) {Entrenamiento\\externo};
      \node[internal,right=of tr] (split) {Crear folds\\internos};
      \node[internal,right=of split] (cand) {Ajustar cada\\pipeline candidato};
      \node[internal,right=of cand] (val) {Validar sin\\reajustar};
      \node[decision,right=of val] (more) {¿Más folds y\\configuraciones?};
      \node[process,right=of more] (pick) {Agregar y elegir\\con la regla fijada};
      \draw[flow] (tr)--(split); \draw[flow] (split)--(cand); \draw[flow] (cand)--(val); \draw[flow] (val)--(more);
      \draw[flow] (more)--node[above] {no} (pick);
      \draw[flow] (more.south) -- ++(0,-.6) -| node[pos=.25,below] {sí} (split.south);
    \end{tikzpicture}
  \end{center}
  \begin{itemize}
    \item Los folds internos respetan la misma unidad, agrupación y temporalidad.
    \item Las métricas se agregan antes de elegir.
    \item No se selecciona por el resultado de un fold aislado.
  \end{itemize}
\end{frame}

% 7
\begin{frame}{El pipeline vive dentro de cada fold}
  \begin{columns}[T]
    \column{.52\textwidth}
    \centering
    \begin{tikzpicture}[node distance=4mm,scale=.75,transform shape]
      \node[data] (tr) {Entrenamiento\\interno};
      \node[internal,below=of tr] (prep) {Ajustar imputación,\\codificación y escala};
      \node[internal,below=of prep] (sel) {Ajustar selección,\\balanceo o PCA};
      \node[internal,below=of sel] (mod) {Ajustar modelo};
      \draw[flow] (tr)--(prep); \draw[flow] (prep)--(sel); \draw[flow] (sel)--(mod);
    \end{tikzpicture}
    \column{.46\textwidth}
    \centering
    \begin{tikzpicture}[node distance=4mm,scale=.75,transform shape]
      \node[data] (va) {Validación\\interna};
      \node[external,below=of va] (apply) {Aplicar objetos\\congelados};
      \node[external,below=of apply] (pred) {Predecir};
      \node[external,below=of pred] (metric) {Medir};
      \draw[flow] (va)--(apply); \draw[flow] (apply)--(pred); \draw[flow] (pred)--(metric);
    \end{tikzpicture}
  \end{columns}
  \begin{alertblock}{Fuga}Preprocesar una sola vez antes de crear los folds permite que validación influya en los parámetros aprendidos.\end{alertblock}
\end{frame}

% 8
\begin{frame}{Predicciones externas y agregación}
  \begin{center}
    \begin{tikzpicture}[node distance=5mm,scale=.78,transform shape]
      \node[external] (f1) {Fold externo 1\\predicciones};
      \node[external,below=of f1] (f2) {Fold externo 2\\predicciones};
      \node[external,below=of f2] (fk) {Fold externo $K$\\predicciones};
      \node[process,right=13mm of f2] (pool) {Integrar predicciones\\fuera de muestra};
      \node[process,right=9mm of pool] (metrics) {Métricas globales,\\por fold y calibración};
      \draw[flow] (f1)--(pool); \draw[flow] (f2)--(pool); \draw[flow] (fk)--(pool); \draw[flow] (pool)--(metrics);
    \end{tikzpicture}
  \end{center}
  \small
  \begin{itemize}
    \item Conservar soporte, composición y resultados por fold.
    \item En CV repetida, no multiplicar artificialmente observaciones independientes.
    \item No iniciar otra selección con resultados externos sin reconocer su reutilización.
  \end{itemize}
\end{frame}

% 9
\begin{frame}{Modelo definitivo y prueba opcional}
  \begin{center}
    \begin{tikzpicture}[node distance=5mm,scale=.61,transform shape]
      \node[data] (dev) {Todo el desarrollo};
      \node[internal,right=of dev] (search) {Nueva búsqueda\\interna};
      \node[process,right=of search] (choose) {Configuración\\elegida};
      \node[process,right=of choose] (fit) {Reajustar pipeline\\completo};
      \node[process,right=of fit] (save) {Persistir modelo,\\transformaciones y umbral};
      \node[optional,right=of save] (test) {Prueba independiente\\una sola vez};
      \draw[flow] (dev)--(search); \draw[flow] (search)--(choose); \draw[flow] (choose)--(fit); \draw[flow] (fit)--(save);
      \draw[flow,dashed,draw=UAPGray] (save)--(test);
    \end{tikzpicture}
  \end{center}
  \begin{itemize}
    \item La evaluación externa estima el procedimiento; no entrega automáticamente el modelo final.
    \item El ajuste definitivo usa todos los datos de desarrollo permitidos.
    \item No se conserva el ajuste del fold externo con mejor puntuación.
  \end{itemize}
\end{frame}

% 10
\begin{frame}{Ejemplo: anticipar demora crítica}
  \scriptsize
  \begin{tabularx}{\textwidth}{>{\bfseries}p{.23\textwidth}X}\toprule
  Nivel & Decisión ilustrativa\\\midrule
  Tarea & demora mayor que 10 minutos\\
  Partición externa & vehículos completos; ningún vehículo aparece en ambos lados\\
  Partición interna & folds por vehículo dentro del entrenamiento externo\\
  Pipelines & logística escalada, SVM con RBF y árbol controlado\\
  Configuraciones & regularización, $C$, $\gamma$, profundidad y umbral\\
  Selección interna & recall mínimo y menor cantidad de falsas alarmas\\
  Evaluación externa & predicciones fuera de muestra por vehículo\\
  Ajuste final & nueva búsqueda interna y reajuste con desarrollo completo\\\bottomrule
  \end{tabularx}
  \idea{La estructura de partición responde a la generalización buscada: vehículos no vistos.}
\end{frame}

% 11
\begin{frame}{Errores que rompen la separación}
  \scriptsize
  \begin{tabularx}{\textwidth}{>{\bfseries}p{.36\textwidth}X}\toprule
  Error & Consecuencia\\\midrule
  preprocesar antes de crear folds & validación influye en la representación\\
  seleccionar con resultados externos & la estimación externa se vuelve optimista\\
  abrir prueba repetidamente & prueba pasa a cumplir función de validación\\
  mezclar grupos o periodos & se evalúa una tarea más fácil que el uso real\\
  reportar el ganador interno & se ignora el sesgo de selección\\
  conservar el mejor fold externo & se elige por azar de esa partición\\\bottomrule
  \end{tabularx}
  \begin{alertblock}{Límite}La anidación organiza el uso de la información; no vuelve válida cualquier decisión experimental.\end{alertblock}
\end{frame}

% 12
\begin{frame}{Síntesis}
  \small
  \begin{itemize}
    \item El pipeline completo es la unidad evaluada.
    \item El ciclo interno selecciona configuraciones.
    \item El ciclo externo estima el procedimiento de selección.
    \item Cada transformación se ajusta solo con el entrenamiento correspondiente.
    \item Las predicciones externas se agregan sin volver a seleccionar.
    \item El modelo definitivo se reajusta después de evaluar el procedimiento.
    \item La prueba es opcional, pero permanece intacta si existe.
  \end{itemize}
  \begin{block}{Simplificación válida}Sin selección adaptativa puede omitirse el ciclo interno, pero no el ajuste local de transformaciones.\end{block}
  \scriptsize Lecturas: capítulo 9, §9.3; capítulo 3, §3.4; capítulo 10, §10.2.7.
\end{frame}

\end{document}
